% ============================================
% DOC-24 (EN)
% English translation — SPEC-935-R355
% ============================================
\section*{\textbf{DOC-24} --- Conservation Technical Assessment}
\markboth{DOC-24 --- Conservation Technical Assessment}{}

\noindent \textit{Document extracted from the forensic file.}

\vspace{0.5em}\noindent\rule{\textwidth}{0.5pt}\vspace{0.5em}

\noindent {\footnotesize\texttt{PUBLIC ARCHIVE $\cdot$ DOCUMENTARY EVIDENCE $\cdot$ CASE 1,260}}\\[-0.2em]

\noindent \textbf{Minas Gerais Public Archive --- Conservation and Restoration Laboratory}\\
\noindent \textbf{Object:} bound notebook, record 1,260.\\
\noindent \textbf{Responsible technician:} Fernanda Lopes Bastos, CRT-04/1128.\\
\noindent \textbf{Request:} assessment of physical integrity and storage recommendation.\\
\noindent\textbf{Date:} August 4, 2026.

\noindent \textbf{1. General description}

\noindent Notebook on paper of variable weight (between 80 and 120\,g/m²), handmade tanned-leather binding, probably local or artisanal production. The general state of conservation is inconsistent with the declared age (dates between 1917 and 1979). The paper does not present the expected signs of oxidation for an object of the period. The acidity of the support is neutral --- unusual for non-industrial paper from the early 20th century.

\noindent \textbf{2. Main finding}

\noindent During the folio analysis, a discrepancy was identified between the thickness of the notebook and the number of pages recorded in the 1979 inventory. The original inventory --- prepared by Dr. Antônio Tavares at the time of the exhumation --- lists approximately 190 filled pages. The current count indicates 208 folios with text.

\noindent There is no evidence that pages were physically added. The binding is intact, without signs of disassembly or reinsertion. Original stitching, original glue. And yet, there are 18 pages more.

\noindent \textbf{3. Analysis of the excess pages}

\noindent The 18 additional pages are distributed throughout the notebook, interleaved between the original folios. In each case, the new page --- or apparently new page --- inserts itself between two pages whose text is continuous and whose handwriting corresponds to the immediately preceding period. The additional page seems to have ``grown'' between the two, as if the notebook had generated new paper to accommodate text that did not exist in the original writing.

\noindent The ink of the additional pages has the same chemical composition as the original ink, but the iron content in the formulation is slightly higher --- indicating more recent oxidation. The relative dating test places the writing of these pages between 2020 and 2026.

\noindent The additional pages were written inside a notebook sealed in a vault since 1979.

\noindent \textbf{4. Content of the additional pages}

\noindent \textit{At the request of the directorate, the content of the 18 pages was not transcribed. I report only that the handwriting does not correspond to the calligraphy of patient 1,260. The handwriting of the additional pages resembles --- preliminarily, without conclusive graphotechnical analysis --- that of Dr. Heitor Oliveira in his last records.}
\noindent \textbf{5. Technical recommendation}

\noindent Storage in a sealed envelope with argon atmosphere. Weekly monitoring of weight and thickness. Prohibition of individual handling. The responsible technician requests --- through this assessment --- not to be summoned for a new analysis of the object.

\noindent \textit{Fernanda Lopes Bastos}\\
\noindent \textit{Conservation and Restoration Technician}\\
\noindent \textit{CRT-04/1128}

\vspace{0.5em}
\noindent\textit{\textbf{↪ Routes:} \rota{DOC-01} \rotasep \rota{LUC-02} \rotasep \rota{DOC-25}}
